\section{Introduction}

Large-scale parallel and distributed systems rely on interconnection networks to coordinate thousands or even millions of processing units. Faults in such systems may corrupt data, interrupt services, and trigger cascading failures. System-level diagnosis offers a classical and effective paradigm for identifying faulty processors through mutual tests among adjacent processors~\cite{pmc67,kavianpour2002comparative}. In the PMC model, a processor tests its neighbors and the collection of test outcomes, called a syndrome, is used to infer the faulty processors in the system.

Most classical diagnosis studies assume deterministic testing behavior and derive diagnosability conditions from specific topological properties~\cite{wang1994diagnosability,zhu2016diagnosis,chang2020conditional,tang2024connectivity}. These results are theoretically important, but they become difficult to use in modern systems where diagnostic observations are noisy and the same test may produce different outcomes under different operating conditions. The probabilistic PMC model addresses this issue by assigning probabilities to test outcomes~\cite{blough1989efficient,sun2023probabilistic}. In this model, the reliability of a test is controlled by parameters such as $\gamma$ and $\beta$, which characterize the probability of observing agreement or disagreement under different node states. However, probabilistic diagnosis is computationally challenging because the true fault set is a latent variable and the likelihood must be evaluated over an exponentially large search space.

Graph neural networks (GNNs) provide a promising tool for scalable diagnosis because interconnection networks are naturally represented as graphs and fault evidence is strongly local. A straightforward solution is to transform probabilistic diagnosis into node-level binary classification and train a GNN to predict whether each processor is faulty. Although effective in simple settings, this classification-oriented formulation has three important limitations. First, a hard label does not reveal how confident the diagnosis is, even when several nodes have ambiguous evidence. Second, the output probability of a neural classifier is often poorly calibrated; a score of $0.9$ does not necessarily mean that the node is faulty with probability $0.9$. Third, existing learning-based formulations usually assume fixed diagnostic parameters during training and testing, while the true reliability of tests may be unknown, heterogeneous, or shifted by workload and environmental changes.

This paper takes a different view. We formulate probabilistic fault diagnosis as an uncertainty-aware posterior inference problem rather than a hard classification problem. Given a graph and multi-round probabilistic syndromes, our goal is to estimate the node-wise fault posterior probability $P(z_u=1\mid \mathbb{S},G)$, quantify the uncertainty of this estimate, and provide a ranked list of suspicious processors for risk-aware inspection. This formulation better matches the probabilistic nature of the PMC model: when the evidence is strong, the model should make confident predictions; when the evidence is conflicting or the testing environment is unreliable, the model should explicitly expose uncertainty instead of producing overconfident labels.

To this end, we propose \tool, an uncertainty-aware probabilistic graph diagnosis framework. \tool combines two complementary sources of diagnostic evidence. The first source is a set of interpretable statistical summaries, such as match ratio, mismatch ratio, and neighborhood disagreement variation. These features provide stable local evidence and support explanation. The second source is the raw multi-round syndrome sequence on each edge, which preserves fine-grained stochastic patterns that may be lost after simple aggregation. A reliability-aware graph inference module then propagates evidence across the topology while learning how much each edge-level observation should be trusted. Finally, a posterior diagnosis head produces calibrated fault probabilities and uncertainty scores. To avoid overfitting to a single diagnostic environment, \tool is trained with randomized probabilistic PMC parameters and evaluated under both seen and unseen parameter settings.

In summary, this paper makes the following contributions:
\begin{itemize}[leftmargin=*, itemsep=1pt, parsep=0pt, topsep=1pt]
    \item We formulate probabilistic PMC diagnosis as a posterior inference problem that estimates calibrated node-wise fault probabilities, uncertainty scores, and top-$k$ suspicious processors instead of only hard fault labels.
    \item We propose \tool, a reliability-aware graph diagnosis framework that fuses statistical diagnostic summaries with raw multi-round syndrome observations for probabilistic evidence propagation.
    \item We introduce parameter-robust training for unknown or varying probabilistic PMC reliability parameters, enabling diagnosis under diagnostic reliability shifts.
    \item We conduct experiments on representative interconnection networks and evaluate not only accuracy, precision, recall, and F1-score, but also calibration quality, ranking effectiveness, and robustness to parameter shifts.
\end{itemize}